% comp s02, from old m16 (cost, data augmentation, approximations, sampler)
\section{What a posterior costs}\label{comp:sec:cost}

Write $Z$ for the design matrix and $\sigma^2\Lambda$ for the prior covariance of the coefficients, with $\Lambda=\mathrm{diag}(\lambda_j^2\gscale^2)$. Each iteration of the Gibbs sampler\index{Markov chain Monte Carlo|textbf}\index{MCMC|see{Markov chain Monte Carlo}} of Chapter~\ref{ch:horseshoe} draws the coefficients from a Gaussian full conditional, which needs one Cholesky factorization\index{Cholesky factorization|textbf} of the $p\times p$ matrix $A=Z^\top Z+\Lambda^{-1}$ at $O(p^3)$, and then draws the scales and the error variance from inverse gamma full conditionals at $O(p^2)$ at most. The cost of a posterior\index{computational cost of a posterior|textbf} is one factorization times the number of iterations: the first factor is known from the size of the model, the second has to be checked.

\subsection{The horseshoe Gibbs sampler}\label{sec:hsgibbs}

The data enter the sampler\index{Gibbs sampler|textbf}\index{Gibbs sampler!for the horseshoe|textbf}\index{horseshoe prior!Gibbs sampler|textbf} only through the cross products $Z^\top Z$, $Z^\top Y$ and $Y^\top Y$, which are formed once at $O(np^2)$. With $A = LL^\top$, a draw of the coefficients is $A^{-1}Z^\top Y + \sigma L^{-\top}\eta$ with $\eta \sim N(0, I_p)$, and the residual sum of squares that the update of $\sigma^2$ needs is $Y^\top Y - 2\beta^\top Z^\top Y + \beta^\top Z^\top Z \beta$. Nothing in an iteration then touches $n$, and for a Gaussian outcome the sampler is as fast for ten thousand units as for two hundred. The simulation sampler timed below forms the residual vector $Y - Z\beta$ in each iteration instead, an $O(np)$ step that is small at $n = 200$.

The scales mix slowly when many coefficients are near zero, so the number of iterations has to be checked and not assumed. Most simulations of this book use 1,500 iterations with the first 500 discarded, the 401(k) analysis 3,000 with the first 1,000 discarded, and the horseshoe fits inside double machine learning in Section~\ref{sec:baipw} 600 sweeps with the first 200 discarded; the checks of Section~\ref{comp:sec:checks} say whether that was enough.

In the simulation of Section~\ref{sec:ricsim}, with $n = 200$ and $p = 51$ columns (the treatment and 50 covariates), 1,500 iterations took 0.31 seconds on one core of an Intel Xeon Gold 6240R at 2.4 GHz, the mean of 20 timed runs (the slowest took 0.33 seconds). One cross validated lasso fit took 0.064 seconds, so the sampler costs about five lasso fits, and the ten lasso fits of a five fold cross fitted\index{cross fitting} estimate took 0.60 seconds, twice the sampler. On the 401(k) data of Chapter~\ref{ch:k401}, with $n = 9{,}915$ and 56 columns (the treatment, an intercept and the 54 terms of a quadratic basis in the nine household characteristics), 3,000 iterations took 0.63 seconds (the mean of five runs, 0.58 to 0.64). On the same data one lasso fit took 0.09 seconds, one random forest fit 1.2 seconds, and the ten forest fits of one cross fitted estimate 12 seconds, about nineteen times the sampler. The lasso is cross validated with five folds, the forest has 200 trees, at least 5 observations per leaf and 4 covariates per split (the settings of the DML comparison in that chapter), and both are fitted on four fifths of the rows, as in one fold. Appendix~\ref{app:code} names the script. For large $p$ the factorization dominates, and when $p > n$ the exact sampler of Bhattacharya, Chakraborty and Mallick (2016) costs $O(n^2 p)$ through the Woodbury identity, orders of magnitude less for $p$ in the thousands and $n$ in the hundreds.

We checked the treatment coefficient, the quantity reported in the tables. In each of 20 data sets of the design of Section~\ref{sec:ricsim} we ran four chains of 1,500 iterations from the sampler's default start, discarded the first 500, and computed the split $\hat R$\index{split R hat@split $\hat R$|textbf} and the effective sample size\index{effective sample size|textbf} from the 4,000 draws left. For the independent horseshoe the largest $\hat R$ was 1.003 and the smallest effective sample size 1,660, and for the coupled horseshoe 1.007 and 1,340. In 10 data sets each of the designs with $p = 150$ and $p = 300$ the largest $\hat R$ was 1.012 and the smallest effective sample size 620, and only one of the 80 simulated data sets, the coupled prior at $p = 150$, had an $\hat R$ above 1.01. In 10 data sets of the design with thirty weak confounders, run the same way, both coupled priors passed in every data set, with the largest $\hat R$ 1.004 and the smallest effective sample size 769, and the independent horseshoe did not: six of the ten had an $\hat R$ above 1.01, the largest 1.020, and eight had fewer than 400 effective draws, the smallest 234. With 234 effective draws the Monte Carlo error of a posterior mean is about a fifteenth of a posterior standard deviation, far less than that design's bias of 0.670, which none of the independent horseshoe's intervals covers; the endpoints of those intervals are less secure. On the 401(k) data, four chains of 3,000 iterations with the first 1,000 discarded gave $\hat R$ of 1.001 for the independent prior and 1.000 for the coupled prior, and an effective sample size equal to the number of draws. Apart from the 600 sweep runs of Section~\ref{sec:baipw}, which we did not check, and the independent horseshoe with thirty weak confounders, the chains behind Chapters~\ref{ch:horseshoe} and \ref{ch:k401} mixed for the treatment coefficient; the check says nothing about the other parameters.

Coefficients that must stay unshrunk, such as an intercept, the treatment coefficient and an estimated propensity, enter $A$ with a prior precision of essentially zero, $10^{-8}$ in the 401(k) analysis of Chapter~\ref{ch:k401}. Unless the outcome and every column are centred, a design without the intercept forces the regression through the origin, and the treatment coefficient, which multiplies the only uncentred column, absorbs the outcome mean. Our first run made this mistake and produced an effect more than twice the correct value, with a deceptively narrow posterior.

\subsection{Data augmentation for binary and count outcomes}

The logistic likelihood is not conjugate to a Gaussian prior, but the Polya Gamma\index{Polya Gamma@P\'olya Gamma} augmentation\index{Polya Gamma@P\'olya Gamma!augmentation|textbf}\index{data augmentation|textbf} of Polson, Scott and Windle (2013) makes it conditionally Gaussian via the identity
\[
\frac{(e^{\zeta})^a}{(1 + e^{\zeta})^b} = 2^{-b}e^{\tilde y\zeta}\int_0^\infty e^{-\omega\zeta^2/2}\,p(\omega \mid b, 0)\,d\omega, \qquad \tilde y = a - b/2,
\]
where $p(\omega\mid b,0)$ is the Polya Gamma density (Appendix~\ref{app:toolkit}). A logistic unit with linear predictor $\zeta$ has $a=y_i$ and $b=1$, so $\tilde y_i=y_i-1/2$ as in Chapter~\ref{ch:horseshoe}. Given the $\omega_i$ the likelihood is a Gaussian regression with weights $\omega_i$ and working response $\tilde y_i/\omega_i$, so the horseshoe step\index{horseshoe prior!in logistic regression|textbf} applies unchanged, and given the linear predictor each $\omega_i$ has a Polya Gamma full conditional that can be sampled exactly. The same device handles negative binomial counts and multinomial outcomes.

Alternating the two draws samples the propensity model of Chapter~\ref{ch:horseshoe}, with the coupled priors that share scales between the treatment and outcome equations, and the coefficients of a binary outcome such as cancer in Chapter~\ref{ch:smoking}. Now $Z^\top\Omega Z+\Lambda^{-1}$, with $\Omega=\mathrm{diag}(\omega_i)$, depends on $\omega$, so $Z^\top\Omega Z$ is formed in every iteration at $O(np^2)$, and unlike the Gaussian sampler this one depends on $n$.

\subsection{When the sampler is too slow}

Two approximations are standard. The Laplace approximation\index{Laplace approximation|textbf} replaces the posterior by a normal centred at the mode with covariance from the Hessian; integrated nested Laplace approximations\index{integrated nested Laplace approximation|textbf}\index{INLA|see{integrated nested Laplace approximation}} (Rue, Martino and Chopin 2009) extend it to latent Gaussian models. Variational Bayes\index{variational inference|textbf} approximates the posterior by the closest member, in Kullback and Leibler divergence, of a tractable family.

For a treatment coefficient with a flat prior and a large sample, the Bernstein von Mises theorem makes the posterior close to normal under regularity conditions, and the Laplace approximation is then good. It is poor for the horseshoe scales, whose posteriors are skewed and bounded below, but those are nuisance parameters that are integrated out. Mean field factorization removes the posterior correlation that confounding adjustment creates between treatment and confounder coefficients, underestimating variance. An effect interval can be too short by a factor that depends on confounding strength, and nothing in the optimization reports it.

Either approximation can supply starting values or screen covariates before a full run, but the reported analyses use neither: every chain starts from the sampler's default values. A tail probability may still be wrong under an approximation that is good at the centre.

When the log posterior can be evaluated and differentiated but the full conditionals have no closed form, Hamiltonian Monte Carlo (HMC)\index{Hamiltonian Monte Carlo|textbf} is the general sampler. It adds a Gaussian momentum, follows the gradient by leapfrog steps and accepts or rejects the end point by a Metropolis step, so the normalizing constant is never needed; the no U turn sampler of Hoffman and Gelman (2014) chooses trajectory lengths during sampling and adapts the step size during warm up. Positive scales are sampled on the log scale with the Jacobian. For the horseshoe, the non centred form $\beta_j=\sigma\gscale\lambda_jz_j$ with $z_j\sim N(0,1)$ can ease funnel geometry, especially when coefficients are weakly informed by the data. Divergent transitions call for a check of posterior geometry; the checks of Section~\ref{comp:sec:checks} still apply. Each leapfrog step needs a gradient that touches all $n$ rows, so the cost is iterations times steps times $np$. 

\paragraph{Running the analyses at scale.} For many outcomes, subgroups or sites, batch the Cholesky factorizations\index{Cholesky factorization} across outcomes that share a design matrix, and start each run at the posterior mean of a related run, which shortens burn in. When one model is fitted to many similarly structured data sets, a generative map\index{generative Bayesian computation!cost|textbf} replaces each run by a forward pass (Section~\ref{comp:sec:algorithm}), with the calibration check of Section~\ref{comp:sec:checks} run once on the trained network.
